% Capítulo 6 (Resultados) — Slide 9 — Dominância por perfil (tabela)
\begin{frame}{Dominância depende do perfil}
  \centering
  \resizebox{0.85\linewidth}{!}{%
  \begin{tabular}{@{}lcrr@{}}
    \toprule
    \textbf{Perfil} & $n$ & \textbf{Dominante} & \textbf{CTI médio (R\$)} \\
    \midrule
    Sparse High Impact & 116 & SA & 584,01 \\
    Unstable Trend & 18* & EOQ & 607,95 \\
    High Vol. Seasonal & 11* & EOQ & 618,79 \\
    \bottomrule
  \end{tabular}}\\[0.3em]
  {\scriptsize *$n<20$: evidência exploratória.}
  \vspace{0.8em}

  \begin{alertblock}{Cuidado com o custo bruto}
    \small O Jornaleiro tem o \highlight{menor CTI bruto} em Sparse High
    Impact (R\$\,284), mas NS\,$=$\,0,55 -- abaixo do limiar de viabilidade
    de 0,70. Não é dominante.
  \end{alertblock}
  \note{
    Esta tabela resume a política dominante por perfil, sempre sob a
    restrição de NS mínimo de 0,70. O SA domina no perfil que concentra 80\%
    das séries; o EOQ domina nos dois perfis menores, com evidência
    exploratória por terem menos de 20 séries. A ressalva mais importante é
    sobre o Jornaleiro: ele tem o menor custo bruto no perfil mais numeroso,
    mas seu nível de serviço de 0,55 está muito abaixo do limiar mínimo --
    por isso é excluído da dominância, mesmo sendo o mais barato. Essa
    ressalva evita a leitura simplista de que custo baixo implica política
    melhor.
  }
\end{frame}
